\chapter{Kết luận và hướng phát triển}
\label{chap:ketluan}

%==============================================================================
\section{Những việc đã làm được}
\label{sec:da-lam}
%==============================================================================

Luận văn đặt ra năm nhiệm vụ ở Mục~\ref{sec:muc-tieu}. Phần này đối chiếu từng
nhiệm vụ với kết quả thực tế, kể cả những chỗ chưa đạt.

\paragraph{MT1 --- Cơ sở toán học của mạng truyền thẳng.}
Đã hoàn thành, với trọng tâm khác tài liệu học sâu tổng quát đúng như dự kiến.
Hai hệ thức truy hồi \eqref{eq:G-truy-hoi} và \eqref{eq:S-truy-hoi} cho Jacobi và
Hessian \emph{theo biến đầu vào} được thiết lập tường minh
(Mệnh đề~\ref{md:jacobian}, \ref{md:hessian}) và trở thành công cụ trung tâm cho
mọi tính toán phần dư về sau. Từ chúng suy ra ngay điều kiện \emph{cần} trên hàm
kích hoạt, mà trường hợp ReLU là dạng cực đoan: hàm mục tiêu phần dư của phương
trình bậc hai \emph{hằng số địa phương}, nên mọi thuật toán gradient đứng yên
tuyệt đối (Hệ quả~\ref{hq:pinn-hong}).

\paragraph{MT2 --- Vi phân tự động và thuật toán tối ưu.}
Đã hoàn thành. Mắt xích quan trọng nhất là Mệnh đề~\ref{md:ad-tai-tao}: cơ chế
thuận-trên-thuận của vi phân tự động \emph{tái tạo đúng} hai hệ thức truy hồi của
Chương~\ref{chap:fnn} mà không cần lập trình tường minh $\sigma''$. Nhờ đó,
Thuật toán~\ref{alg:dao-ham} đóng vai trò đặc tả toán học, còn cài đặt thực tế
chỉ là vài dòng gọi \texttt{torch.autograd.grad} với cờ
\texttt{create\_graph=True}. Chiến lược hai giai đoạn Adam~$\to$~L-BFGS được biện
minh bằng ba lập luận cụ thể --- tính tất định của gradient, sàn nhiễu
(Mệnh đề~\ref{md:san-nhieu}), và điều kiện xấu --- chứ không bằng kinh nghiệm
thực hành.

\paragraph{MT3 --- Xây dựng PINNs một cách chặt chẽ.}
Đã hoàn thành, và đây là chỗ luận văn bổ sung được một mắt xích mà tài liệu
thường bỏ qua: \emph{vì sao cực tiểu hoá phần dư lại là việc đáng làm}.
Mệnh đề~\ref{md:chan-on-dinh} trả lời bằng một chặn ổn định, và
Mệnh đề~\ref{md:chan-suy-bien} chỉ ra hằng số của chặn ấy suy biến theo $1/\nu$.
Đóng góp tổng hợp của chương là Bảng~\ref{tab:ba-nguyen-nhan}: ba nguyên nhân độc
lập, nằm ở ba tầng khác nhau, cùng làm PINNs suy giảm trên đúng lớp bài toán có
thang độ dài nhỏ.

\paragraph{MT4 --- Cài đặt và thực nghiệm trên PyTorch.}
Đã hoàn thành. Toàn bộ cài đặt dùng PyTorch thuần, không dùng thư viện PINNs đóng
gói sẵn, và được công bố kèm luận văn (Phụ lục~\ref{app:code}) để mọi bảng số và
mọi hình trong Chương~\ref{chap:thucnghiem} đều tái lập được bằng một lệnh.

\paragraph{MT5 --- Khảo sát độ nhạy và phân tích hạn chế.}
Đã hoàn thành ở mức tương xứng với ngân sách tính toán, nhưng đây cũng là nhiệm
vụ còn nhiều dư địa nhất; Mục~\ref{sec:hanche} đã nêu thẳng tám hạn chế của chính
thiết kế thực nghiệm.

%==============================================================================
\section{Đóng góp của luận văn}
\label{sec:dong-gop-cuoi}
%==============================================================================

Trong phạm vi đã nêu ở Mục~\ref{subsec:pham-vi}, luận văn có ba đóng góp.

\emph{Thứ nhất}, một trình bày mạch lạc bằng tiếng Việt về PINNs, đi từ cơ sở
toán học của mạng truyền thẳng tới công thức đầy đủ của mô hình thời gian liên
tục, với notation nhất quán xuyên suốt và mọi khẳng định đều truy được về một
mệnh đề cụ thể.

\emph{Thứ hai}, một cài đặt PyTorch độc lập kèm kết quả thực nghiệm định lượng,
tái lập được, có thể dùng làm điểm khởi đầu cho các nghiên cứu tiếp theo.

\emph{Thứ ba}, một phân tích đối chiếu có hệ thống giữa lý thuyết và số liệu
(Bảng~\ref{tab:doi-chieu}). Phần đối chiếu này cho một kết luận phương pháp luận
đáng chú ý: những hạng mục khớp trọn vẹn đều tương ứng với một mệnh đề đã được
chứng minh, còn mọi hạng mục phải điều chỉnh đều tương ứng với chỗ suy luận tiệm
cận hoặc ghép hai kết quả rời rạc.

Bên cạnh ba đóng góp đã dự kiến, quá trình thực nghiệm sinh thêm hai kết quả
không nằm trong kế hoạch ban đầu, và cả hai đều xuất phát từ \emph{kết quả tiêu
cực}:

\begin{itemize}[leftmargin=*]
\item \textbf{Một điều chỉnh phát biểu lý thuyết.} Mục~\ref{subsec:he-qua-pinn}
      kết luận rằng bài toán đòi hỏi độ chính xác cao nhất ở đúng những tần số mà
      mô hình học chậm nhất. Thí nghiệm~\ref{tn:pho} cho thấy phát biểu này
      không đúng phổ quát: với toán tử elliptic thuần tuý, hàm mục tiêu phần dư
      \emph{ưu tiên} tần số cao và chính tần số thấp mới bị bỏ rơi. Cách phát
      biểu đã được siết lại ở Nhận xét~\ref{nx:hieu-chinh}.
\item \textbf{Một chế độ hỏng chứng minh được của quy tắc cân bằng trọng số.}
      Thuật toán~\ref{alg:annealing} làm giảm độ chính xác trên cả hai bài toán
      khảo sát. Truy nguyên dẫn tới Bổ đề~\ref{bd:phan-hoi-duong}: quy tắc
      \eqref{eq:lambda-hat} có một vòng phản hồi dương với tốc độ phân kỳ
      $\hat\lambda_i \gtrsim J_i^{-1/2}$ và không có cơ chế tự hãm. Đây là kết
      quả lý thuyết nảy sinh \emph{từ} thực nghiệm chứ không phải được kiểm chứng
      bởi thực nghiệm.
\end{itemize}

%==============================================================================
\section{Hạn chế}
\label{sec:han-che-cuoi}
%==============================================================================

Ngoài tám hạn chế về thiết kế thực nghiệm đã nêu ở Mục~\ref{sec:hanche}, cần nói
rõ ba giới hạn về phạm vi của toàn bộ luận văn.

Thứ nhất, luận văn \emph{không} chứng minh định lý mới về hội tụ hay ước lượng
sai số cho PINNs. Các kết quả lớn của lĩnh vực chỉ được trích dẫn cùng diễn giải;
những chứng minh được trình bày đầy đủ đều sơ cấp, và giá trị của chúng nằm ở chỗ
biến những khẳng định mơ hồ thành phát biểu có giả thiết rõ ràng.

Thứ hai, phạm vi bài toán hẹp: chỉ mô hình thời gian liên tục, chỉ mạng truyền
thẳng kết nối đầy đủ, chỉ bài toán thuận, và chỉ các phương trình một chiều theo
không gian. Lợi thế lý thuyết quan trọng nhất của mạng nơ-ron --- bậc xấp xỉ
$\mathcal{O}(m^{-1/2})$ không phụ thuộc số chiều ở Định lý~\ref{dl:barron} ---
vì vậy \emph{không} được kiểm chứng bằng thực nghiệm nào trong luận văn.

Thứ ba, các chặn ổn định ở Mệnh đề~\ref{md:chan-on-dinh} và
Mệnh đề~\ref{md:chan-suy-bien} được thiết lập cho bài toán tuyến tính một chiều
với điều kiện biên thoả chính xác. Với ràng buộc mềm và với toán tử phi tuyến,
chúng chỉ đóng vai trò cơ sở khái niệm chứ không phải chặn hậu nghiệm áp dụng
được trực tiếp.

%==============================================================================
\section{Hướng phát triển}
\label{sec:huong-phat-trien}
%==============================================================================

Bảng~\ref{tab:ba-nguyen-nhan} gợi ý một nguyên tắc sắp xếp các hướng tiếp theo:
vì ba nguyên nhân nằm ở ba tầng khác nhau, mỗi hướng nên được đánh giá theo
\emph{tầng mà nó tác động tới}, và không hướng nào thay thế được hai tầng còn
lại.

\paragraph{Tầng bài toán.} Phân rã miền và hành quân theo thời gian
(Bảng~\ref{tab:remedies}) thay một bài toán khó bằng một dãy bài toán dễ hơn, tức
tác động trực tiếp lên hằng số ổn định. Đây là hướng đáng thử trước với phương
trình Burgers độ nhớt nhỏ, nơi Mục~\ref{sec:tn5} cho thấy sai số tập trung vào
một dải hẹp quanh lớp trong.

\paragraph{Tầng tối ưu.} Bổ đề~\ref{bd:phan-hoi-duong} không chỉ ra rằng cân bằng
trọng số là sai, mà chỉ ra \emph{cơ chế} hỏng. Hai sửa đổi đi thẳng từ cơ chế
ấy đáng được khảo sát: đặt chặn trên cho $\lambda_i$, hoặc thay mẫu số của
\eqref{eq:lambda-hat} bằng một đại lượng không triệt tiêu cùng $J_i$ --- chẳng
hạn vết của nhân tiếp tuyến.

\paragraph{Tầng xấp xỉ.} Thí nghiệm~\ref{tn:pho} cho thấy nhúng đặc trưng Fourier
cải thiện $161$ lần nhưng chỉ \emph{dịch chuyển} dải tần hiệu dụng chứ không mở
rộng nó, và việc chọn thang $s$ đòi hỏi biết trước thang độ dài của nghiệm chưa
biết. Một quy tắc chọn $s$ dựa trên dữ liệu --- ví dụ ước lượng từ phổ của phần
dư trong vài trăm vòng lặp đầu --- là hướng cụ thể và khả thi.

\paragraph{Củng cố phương pháp luận thực nghiệm.} Ba việc rẻ và có giá trị ngay:
lặp mỗi cấu hình trên $5$--$10$ hạt giống rồi báo cáo trung vị cùng khoảng phân
vị (Mục~\ref{sec:hanche}); áp dụng lấy mẫu thích nghi RAR cho bài toán Burgers;
và mở rộng phép kiểm chứng chặn ổn định ở Mục~\ref{sec:tn6} sang bài toán đối
lưu--khuếch tán để đo trực tiếp sự suy biến theo $1/\nu$ mà
Mệnh đề~\ref{md:chan-suy-bien} dự báo.

\paragraph{Mở rộng phạm vi.} Cuối cùng, hai mở rộng nằm ngoài phạm vi luận văn
nhưng nối liền với nó: bài toán ngược, nơi khung hàm mất mát \eqref{eq:total-loss}
đã sẵn sàng và chỉ cần thêm $\bzeta$ vào tập biến tối ưu
(Mục~\ref{subsec:thuan-nguoc}); và bài toán nhiều chiều, nơi lập luận chống lời
nguyền số chiều ở Nhận xét~\ref{nx:loi-nguyen} lần đầu tiên trở nên kiểm chứng
được.
